\section{Step 22 Prepare fresh controls for this block}\label{rgb-s22}

Copy the preserved controls into \texttt{w4}, \texttt{w5} and \texttt{w6}. These working registers will supply controls to \texttt{vshf} and then receive its output. Copying each iteration prevents a prior block's component values from becoming the next block's selectors. A move copies all 128 bits; it does not shuffle them.

\Needspace{13\baselineskip}
\subsection{Write or inspect}
Replace the temporary branch and nop in \texttt{vector\_pack\_copy}. Preserve surrounding labels and the provided entry and return.

\Needspace{6\baselineskip}
\MARSCodeLabel{MIPS assembly}
\begin{minted}[samepage=true]{mars-idea}
    move.v $w4, $w16
    move.v $w5, $w17
    move.v $w6, $w18
\end{minted}
Use \texttt{s22-1} (count 8) from the companion input recipes. Pause at \texttt{pack\_ready} before comparing the construction state.

\subsection{Expected observation}
The working registers now equal their corresponding preserved control vectors.

\Needspace{11\baselineskip}
\subsection{Check your understanding}
Why copy rather than repeatedly use a prior shuffle result as controls?

\begin{itemize}
  \item \textbf{A} The destination is overwritten by each shuffle.
  \item \textbf{B} The preserved table contains pixel addresses.
  \item \textbf{C} A vector move decrements count.
\end{itemize}
